% ============================================================
%  Chapter 3 – System Design and Implementation
% ============================================================
% Ensure English caption labels even inside a Vietnamese template
\renewcommand{\tablename}{Table}
\renewcommand{\figurename}{Figure}

\chapter{System Design and Implementation}

% ------------------------------------------------------------
\section{System Architecture Overview}
% ------------------------------------------------------------

The traffic signal control system is built as a closed-loop interaction between
a Traffic Controller and the SUMO simulation environment, as illustrated in
Figure~\ref{fig:architecture}. At each decision step $t$, each agent $i$
observes its local state $s_i(t)$ at its intersection, passes that state
through a shared policy $\pi$ to select an action $a_i(t)$, and executes the
action inside SUMO. The environment responds with a reward $r_i(t)$ computed
by one of the three reward functions under investigation (queue, pressure, or
wait-clip), and transitions to a new state $s_i(t+1)$. The transition tuple
$(s, a, r, s')$ is collected to update the policy via one of the three
algorithms PPO, DQN, or DDQN. The new state then becomes the input for the
next decision step, closing the learning loop.

In this multi-agent architecture all intersections share a single neural
network (parameter sharing), while each intersection acts independently based
on its own local observation. The interaction between Python and SUMO is
performed through the \texttt{libsumo} interface — a directly embedded variant
of TraCI — which substantially reduces communication latency compared with
standard TraCI.

\begin{figure}[htbp]
    \centering
    \includegraphics[width=0.85\textwidth]{{Tracffic controller}.png}
    \caption{Overall architecture of the reinforcement-learning-based traffic
    signal control system.}
    \label{fig:architecture}
\end{figure}

% ------------------------------------------------------------
\section{Experimental Scenarios}
% ------------------------------------------------------------

Experiments are conducted on two road networks extracted from the TAPAS
Cologne dataset~\cite{tapascologne}, whose characteristics are summarised in
Table~\ref{tab:scenarios}.

\begin{table}[htbp]
    \caption{Parameters of the two experimental scenarios.}
    \label{tab:scenarios}
    \centering
    \begin{tabular}{lcc}
        \toprule
        \textbf{Attribute} & \textbf{Cologne-3} & \textbf{Cologne-8} \\
        \midrule
        Number of TLS       & 3   & 8   \\
        Number of nodes     & 29  & 78  \\
        Number of edges     & 48  & 149 \\
        Number of vehicles  & 2{,}856 & 2{,}046 \\
        \bottomrule
    \end{tabular}
\end{table}

Cologne-3 is a single-axis arterial network with three signalised
intersections, well-suited for observing cascading congestion propagation.
Cologne-8 is a $2\times 2$ grid network with eight signalised intersections
that requires more complex coordination among agents.

\begin{figure}[htbp]
    \centering
    \begin{subfigure}[b]{0.45\textwidth}
        \centering
        \includegraphics[width=\textwidth]{cologne3.png}
        \caption{Cologne-3 (3 signalised intersections).}
        \label{fig:map_c3}
    \end{subfigure}
    \hfill
    \begin{subfigure}[b]{0.45\textwidth}
        \centering
        \includegraphics[width=\textwidth]{cologne8.png}
        \caption{Cologne-8 (8 signalised intersections).}
        \label{fig:map_c8}
    \end{subfigure}
    \caption{Road networks used in the experiments, visualised in SUMO-GUI.
    Signalised intersections are highlighted with traffic-light indicators.}
    \label{fig:maps}
\end{figure}

\noindent\textbf{Training / evaluation split.}\;
Training traffic is synthetically generated with vehicle density matching the
actual count in the original dataset, but with origin–destination
distributions randomised uniformly across the network. This forces agents to
learn a well-generalised policy rather than memorising fixed traffic flows.
Final evaluation is performed on the real TAPAS Cologne traffic using the
original OD matrix, ensuring the authenticity of the reported results.

Vehicles in the simulation use the Krauss car-following model — the default
microscopic model in SUMO~\cite{lopez2018sumo} — with the parameters listed in
Table~\ref{tab:krauss}. The Krauss model simulates safe-driving behaviour with
a stochastic component governed by $\sigma$, capturing inter-driver reaction
variability found in practice.

\begin{table}[htbp]
    \caption{Krauss vehicle model parameters.}
    \label{tab:krauss}
    \centering
    \begin{tabular}{lc}
        \toprule
        \textbf{Parameter} & \textbf{Value} \\
        \midrule
        Car-following model         & Krauss \\
        Maximum acceleration        & $3.28\ \text{m/s}^2$ \\
        Deceleration                & $7.50\ \text{m/s}^2$ \\
        Randomness factor ($\sigma$)& $0.1$ \\
        Reaction time ($\tau$)      & $1.1\ \text{s}$ \\
        Vehicle length              & $4.0\ \text{m}$ \\
        Minimum gap                 & $0.38\ \text{m}$ \\
        \bottomrule
    \end{tabular}
\end{table}

% ------------------------------------------------------------
\section{Simulation Setup}
% ------------------------------------------------------------

All experiments are run on SUMO version~1.26~\cite{lopez2018sumo}. Each episode
spans 3{,}600 simulated seconds (equivalent to one hour of traffic), which is
sufficient for agents to experience multiple traffic cycles including peak-hour
conditions. The simulation time step is 1~s; the agent decision step follows
the signal-phase cycle, meaning agents issue decisions only at phase-transition
moments rather than every second. Total training length is 500{,}000 timesteps
per combination (algorithm $\times$ reward function $\times$ scenario), with
checkpoints saved every 25{,}000 timesteps. The algorithms are implemented in
Python using the Stable-Baselines3 library~\cite{stable_baselines3}, interacting
with SUMO via the \texttt{libsumo} interface to minimise communication overhead.

% ------------------------------------------------------------
\section{Problem Formulation}
% ------------------------------------------------------------

Building on the MDP and Dec-POMDP theoretical foundations presented in
Chapter~2, this section specifies the state space, action space, and reward
function for the traffic signal control problem in the two Cologne scenarios.

% - - - - - - - - - - - - - - - - - - - - - - - - - - - - - -
\subsection{State Representation}
% - - - - - - - - - - - - - - - - - - - - - - - - - - - - - -

The state space must supply the neural network with the most essential
information about traffic conditions at each intersection, and must have a
fixed dimensionality to support parameter sharing across agents. The study
adopts a minimal 7-dimensional (7D) state space, summarised in
Table~\ref{tab:state}.

\begin{table}[htbp]
    \caption{Features of the 7-dimensional state vector.}
    \label{tab:state}
    \centering
    \begin{tabular}{llcc}
        \toprule
        \textbf{Feature} & \textbf{Symbol} & \textbf{Dim.}
            & \textbf{Normalisation} \\
        \midrule
        Green time of 4 phases
            & $g_1,g_2,g_3,g_4$  & 4
            & Divide by 3 (max.\ phases) \\
        Current phase index
            & $P(t)$             & 1
            & Divide by $N_\mathrm{max}$ (max.\ phases) \\
        Remaining time of current phase
            & $T_\mathrm{rem}(t)$& 1
            & Divide by $120.0\ \text{s}$ \\
        Total queue length
            & $Q(t)$             & 1
            & Divide by 50 vehicles \\
        \bottomrule
    \end{tabular}
\end{table}

At each time step $t$, the local observation of agent $i$ is
$s_i(t) = \langle G(t),\, P(t),\, T_\mathrm{rem}(t),\, Q(t) \rangle$.
The four-dimensional component $G(t) = (g_1, g_2, g_3, g_4)$ encodes the
programmed green-time plan for each signal phase, normalised by dividing by
$3$ — the maximum number of phases across all Cologne scenarios.
$P(t)$ is the current phase index, normalised by dividing by $N_\mathrm{max} = 3$ — the maximum number of phases across all Cologne scenarios — to bring its value into $[0, 1]$ on the same scale as the other features. Although $P(t)$ is inherently a categorical variable, one-hot encoding was not used in order to preserve the fixed 7D vector dimensionality required by the parameter-sharing architecture. $T_\mathrm{rem}(t)$ is the remaining green time of
the current phase, divided by 120~s — the maximum permitted green time in the
simulation configuration. $Q(t)$ is the total number of vehicles stopped at
the intersection, divided by 50 — an empirically estimated maximum queue
capacity on the Cologne map. All continuous features are normalised to
$[0, 1]$ before being fed into the neural network, promoting stable training.

The 7D design minimises noise and focuses the model on the most informative
signals. Its main limitation is that each agent observes only local
information; the absence of neighbouring-intersection state may constrain
network-wide coordination.

% - - - - - - - - - - - - - - - - - - - - - - - - - - - - - -
\subsection{Action Formulation}
% - - - - - - - - - - - - - - - - - - - - - - - - - - - - - -

\noindent\textbf{Action space.}\;
The action space at intersection $i$ is defined as
$\mathrm{Discrete}(N_i)$, where $N_i$ is the number of non-conflicting green
phases read directly from the SUMO configuration file. The study adopts the
phase-selection mechanism from the RESCO benchmark~\cite{ault2021resco},
whereby the agent selects which phase to activate rather than adjusting
per-phase durations.

The two main paradigms in RL-based traffic signal control are phase selection
(discrete) and duration adjustment (continuous)~\cite{tan2021hybrid}. Phase
selection uses a discrete action space whose size equals the number of phases,
whereas duration adjustment requires a continuous or higher-dimensional
discrete space to encode the duration of each phase. Research has shown that
agent convergence speed depends considerably on action-space size~\cite{tan2021hybrid},
and phase selection — with its compact discrete space — generally leads to more
stable convergence, particularly in a multi-agent setting~\cite{ault2021resco}.
Moreover, this approach reflects real-world operation where signal phases are
fixed according to traffic-safety standards. However, because each phase
duration is fixed, the agent cannot extend the green time for a direction
experiencing a traffic surge — a flexibility that duration adjustment can
provide, albeit at the cost of a far more complex optimisation problem.

\noindent\textbf{Phase-transition mechanism.}\;
When an agent selects a phase different from the current one, the system
automatically inserts a 5-second yellow transition phase to ensure traffic
safety before activating the new phase. In the multi-agent environment, all
intersections receive actions in parallel from the shared policy, but each
intersection applies its action only at its own decision instant, preserving
independence and flexibility among intersections.

% - - - - - - - - - - - - - - - - - - - - - - - - - - - - - -
\subsection{Reward Functions}
% - - - - - - - - - - - - - - - - - - - - - - - - - - - - - -

The study investigates three reward functions proposed by Nandakumar
et~al.~\cite{nandakumar2025}, each corresponding to a different way of
quantifying congestion. Table~\ref{tab:rewards} summarises the formulae and
value ranges of all three functions.

\begin{table}[htbp]
    \caption{Summary of the three reward functions investigated.}
    \label{tab:rewards}
    \centering
    \begin{tabular}{lcc}
        \toprule
        \textbf{Function} & \textbf{Compact formula} & \textbf{Range} \\
        \midrule
        Queue &
            $\alpha_\mathrm{red}\cdot R_\mathrm{red}
             + \alpha_\mathrm{abs}\cdot R_\mathrm{abs}$ &
            $(-\infty, 1]$ \\
        Pressure &
            $\mathrm{clip}\!\left(
              \dfrac{-(Q_\mathrm{in}-Q_\mathrm{out})}{N_\mathrm{in}\cdot C},
              -1,\, 0\right)$ &
            $[-1, 0]$ \\
        Wait-clip &
            $\mathrm{clip}\!\left(
              \dfrac{-\sum_{v\in V} W_v}{100},
              -5,\, 0\right)$ &
            $[-5, 0]$ \\
        \bottomrule
    \end{tabular}
\end{table}

\noindent\textbf{Queue reward.}\;
The queue-based reward function combines two components: a reward for reducing
the queue relative to the previous step ($R_\mathrm{red}$) and a penalty on
the current absolute queue length ($R_\mathrm{abs}$), following the design
of~\cite{nandakumar2025}:
\begin{equation}
    R_\mathrm{queue} = \alpha_\mathrm{red}\cdot R_\mathrm{red}
                     + \alpha_\mathrm{abs}\cdot R_\mathrm{abs},
\end{equation}
\begin{equation}
    R_\mathrm{red} =
        \frac{\displaystyle\sum_{d\in\mathcal{D}}
              \bigl(Q_d^{t-1} - Q_d^{t}\bigr)}
             {N_\mathrm{act}\cdot C},
    \qquad
    R_\mathrm{abs} =
        \frac{\displaystyle -\sum_{d\in\mathcal{D}} Q_d^{t}}
             {N_\mathrm{act}\cdot C},
\end{equation}
where the notation is defined as follows:
$\mathcal{D} = \{N, S, E, W\}$ is the set of four approach directions
(North, South, East, West);
$Q_d^{t}$ is the maximum number of stopped vehicles (speed equal to zero)
across all lanes of direction $d$ at time step $t$;
$Q_d^{t-1}$ is the corresponding value at the previous step;
$N_\mathrm{act}$ is the number of active approach directions at the
intersection;
$C = 15$ is a normalisation constant empirically estimating the maximum number
of vehicles that can queue on one approach direction;
$\alpha_\mathrm{red} = 0.6$ and $\alpha_\mathrm{abs} = 0.4$ are the weights of
the two components.
The component $R_\mathrm{red}$ rewards every vehicle cleared from the queue
relative to the previous step, providing a dense training signal at each
decision step. The component $R_\mathrm{abs}$ penalises the absolute queue
magnitude, preventing the agent from accumulating a large queue even when the
queue is nominally decreasing.

\noindent\textbf{Pressure reward.}\;
The pressure reward originates from the Max-Pressure control theory proposed
by Varaiya~\cite{maxpressure2013}, was introduced into RL-based reward design
by Wei et~al.\ in PressLight~\cite{presslight2019}, and is applied in the
present experimental setting by Nandakumar et~al.~\cite{nandakumar2025}:
\begin{equation}
    R_\mathrm{pressure} =
        \mathrm{clip}\!\left(
          \frac{-(Q_\mathrm{in} - Q_\mathrm{out})}
               {N_\mathrm{in}\cdot C},\;
          -1,\; 0
        \right),
\end{equation}
where $Q_\mathrm{in}$ is the total number of stopped vehicles across all
incoming lanes; $Q_\mathrm{out}$ is the corresponding count on outgoing lanes;
$N_\mathrm{in}$ is the number of incoming lanes; and $C = 15$ is the
normalisation constant defined above. High pressure (many vehicles entering,
few exiting) produces a large negative reward, incentivising the agent to
prioritise clearance of congested directions. The hard clamp to $[-1, 0]$
stabilises gradients during training. However, the function aggregates
information across all incoming and outgoing lanes without distinguishing which
direction is locally congested within the intersection.

\noindent\textbf{Wait-clip reward.}\;
The wait-clip reward was developed by Nandakumar et~al.\ based on the original
design in the RESCO benchmark~\cite{nandakumar2025,ault2021resco}:
\begin{equation}
    R_\mathrm{wait} =
        \mathrm{clip}\!\left(
          \frac{-\displaystyle\sum_{v\in V} W_v}{100},\;
          -5,\; 0
        \right),
\end{equation}
where $V$ is the set of all vehicles within the control zone of the
intersection; $W_v$ is the cumulative waiting time in seconds of vehicle $v$,
i.e.\ the total time the vehicle has been stationary since entering the
observation zone; the constant 100 is used for scaling to a magnitude
comparable with the other two functions; and the clip threshold $-5$ bounds the
minimum reward to prevent excessively large gradients that could destabilise
training. This function most directly reflects road-user experience and aligns
most closely with the practical objective. However, because $W_v$ accumulates
significantly only after many consecutive stationary steps, the reward signal
appears later than those of the queue and pressure functions, which may slow
learning in early training.

% ------------------------------------------------------------
\section{Model Configuration}
% ------------------------------------------------------------

% - - - - - - - - - - - - - - - - - - - - - - - - - - - - - -
\subsection{Neural Network Architecture}
% - - - - - - - - - - - - - - - - - - - - - - - - - - - - - -

All three algorithms use a two-hidden-layer multi-layer perceptron with
64 neurons per layer and Tanh activation. The input is the 7-dimensional
observation vector; the output is a probability distribution over actions for
PPO, or the Q-value for each action for DQN/DDQN. This architecture is
consistent with those reported in MARL traffic signal control studies of
comparable scale~\cite{ault2021resco} and is well matched to the low-dimensional
input of the 7D state space.

% - - - - - - - - - - - - - - - - - - - - - - - - - - - - - -
\subsection{Hyperparameters}
% - - - - - - - - - - - - - - - - - - - - - - - - - - - - - -

Tables~\ref{tab:ppo_hyper} and~\ref{tab:dqn_hyper} list the hyperparameters of
PPO and DQN/DDQN, respectively. The values are drawn from the Stable-Baselines3
defaults~\cite{stable_baselines3} and tuned through preliminary pilot experiments.

\begin{table}[htbp]
    \caption{PPO hyperparameters.}
    \label{tab:ppo_hyper}
    \centering
    \begin{tabular}{lcp{6cm}}
        \toprule
        \textbf{Parameter} & \textbf{Value} & \textbf{Description} \\
        \midrule
        Learning rate      & $3\times10^{-4}$ & Adam optimiser step size \\
        Rollout steps      & 512  & Steps collected before each update \\
        Batch size         & 128  & Mini-batch size during training \\
        Epochs             & 10   & Passes over each collected rollout \\
        Discount factor $\gamma$ & 0.99 & Weight for future rewards \\
        GAE lambda         & 0.95 & Advantage estimation smoothing \\
        Clip range $\varepsilon$ & 0.2 & Policy-update ratio bound \\
        \bottomrule
    \end{tabular}
\end{table}

Under the Independent multi-agent architecture, each simulation step produces
$N$ transition tuples in parallel — one per intersection. After each 512-step
rollout, the total accumulated samples equal $512\times N$. Over 500{,}000
training timesteps, the number of gradient updates is therefore
$500{,}000 / (512 \times N)$, corresponding to approximately 325 updates for
Cologne-3 ($N = 3$) and 122 updates for Cologne-8 ($N = 8$). The clip
range $\varepsilon = 0.2$ bounds the policy-change ratio at each update step,
ensuring stable training in accordance with the PPO mechanism.

\begin{table}[htbp]
    \caption{DQN and DDQN hyperparameters.}
    \label{tab:dqn_hyper}
    \centering
    \begin{tabular}{lcp{5.5cm}}
        \toprule
        \textbf{Parameter} & \textbf{Value} & \textbf{Description} \\
        \midrule
        Learning rate      & $1\times10^{-4}$ & Adam optimiser step size \\
        Replay buffer size & $500{,}000\times N$
            & 1{,}500{,}000 (C3); 4{,}000{,}000 (C8) \\
        Learning start     & 5{,}000 steps & Steps before first update \\
        Batch size         & 128  & Samples drawn from buffer per update \\
        Discount factor $\gamma$ & 0.99 & Weight for future rewards \\
        Training frequency & 4 steps & Gradient update every 4 sim.\ steps \\
        Target-network update & 1{,}000 steps & Target-network sync interval \\
        Final $\varepsilon$ & 0.05 & Minimum exploration rate (5\,\%) \\
        $\varepsilon$-decay period & 25\,\% of timesteps
            & Linear decay schedule \\
        \bottomrule
    \end{tabular}
\end{table}

The learning rate $1\times10^{-4}$ for DQN/DDQN is lower than the
$3\times10^{-4}$ used for PPO. This reflects a fundamental difference between
the two algorithm families: PPO directly updates the policy on freshly
collected rollout data, whereas DQN/DDQN optimises the Q-function through
repeated random sampling from a replay buffer with bootstrapped targets — a
procedure that is inherently prone to oscillations if the learning rate is too
high~\cite{mnih2015dqn}. The replay buffer is scaled by the number of agents $N$
to ensure a consistent ratio between stored data and update frequency across
the two scenarios. DDQN differs from DQN solely in how the target value is
computed: the action is selected by the online network but evaluated by the
target network, reducing the Q-value overestimation bias described in
Chapter~2.

% - - - - - - - - - - - - - - - - - - - - - - - - - - - - - -
\subsection{Experimental Setup}
% - - - - - - - - - - - - - - - - - - - - - - - - - - - - - -

Each combination (algorithm $\times$ reward function $\times$ scenario) is
trained with 5 independent random seeds, yielding a total of 90 models
$(3\ \text{algorithms} \times 3\ \text{reward functions} \times
2\ \text{scenarios} \times 5\ \text{seeds})$.
Using multiple seeds allows assessment of per-configuration stability and
enables statistically meaningful comparisons across algorithms.

% ------------------------------------------------------------
\section{Baseline Controllers}
% ------------------------------------------------------------

To assess the effectiveness of the reinforcement learning algorithms, the study
compares against two baseline controllers representing two distinct control
paradigms.

% - - - - - - - - - - - - - - - - - - - - - - - - - - - - - -
\subsection{Fixed-Time Control}
% - - - - - - - - - - - - - - - - - - - - - - - - - - - - - -

The Fixed-time controller uses cycle lengths and phase durations defined in the
SUMO configuration file, read directly from the \texttt{duration} attribute of
each \texttt{<phase>} element in the \texttt{.net.xml} file. The controller
repeats this cycle throughout the simulation without adjusting to actual traffic
conditions. This approach represents the traditional paradigm with the
advantage of simplicity and predictability, but is entirely unable to adapt to
traffic fluctuations.

% - - - - - - - - - - - - - - - - - - - - - - - - - - - - - -
\subsection{Webster's Method}
% - - - - - - - - - - - - - - - - - - - - - - - - - - - - - -

Webster's method applies the signal-cycle optimisation formula proposed by
Webster~\cite{webster1958}, in which the optimal cycle length $C_o$ is derived
from the actual-to-saturation flow ratio:
\begin{equation}
    C_o = \frac{1.5\,L + 5}{1 - \sum_i Y_i},
\end{equation}
where $L$ is the total lost time per cycle and $Y_i$ is the flow-to-saturation
ratio of phase $i$. Green time is allocated proportionally to the corresponding
$Y_i$.

In the experimental implementation, Webster operates in dynamic mode: flow is
estimated by a 15-minute moving average and the cycle is recomputed every
145~s. The reference saturation flow is set to the standard 1{,}800~veh/h/lane
($0.5$~veh/s/lane). Cycle length is clamped to $[30,\,120]$~s and the minimum
green time per phase is 5~s. The 145-second update interval is chosen so that
Webster can adapt to flow variations while maintaining cycle stability, and also
ensures a fair comparison with the reinforcement learning algorithms.




% ============================================================
%  Bibliography entries to add to your .bib file
% ============================================================
%
% @inproceedings{lopez2018sumo,
%   author    = {P. A. Lopez and M. Behrisch and L. Bieker-Walz and
%                J. Erdmann and Y.-P. Fl{\"o}tter{\"o}d and R. Hilbrich
%                and L. L{\"u}cken and J. Rummel and P. Wagner and E. Wi{\ss}ner},
%   title     = {Microscopic Traffic Simulation using {SUMO}},
%   booktitle = {Proc.\ IEEE Intelligent Transportation Systems Conference
%                (ITSC)},
%   year      = {2018},
%   doi       = {10.1109/ITSC.2018.8569938}
% }
%
% @misc{tapascologne,
%   title  = {{TAPAS} {Cologne} Scenario},
%   author = {{SUMO Development Team}},
%   howpublished = {\url{https://sumo.dlr.de/docs/Data/Scenarios/TAPASCologne.html}},
%   note   = {Accessed: 2025}
% }
%
% @inproceedings{ault2021resco,
%   author    = {J. Ault and G. Sharon},
%   title     = {Reinforcement Learning Benchmarks for Traffic Signal Control},
%   booktitle = {NeurIPS Datasets and Benchmarks Track},
%   year      = {2021}
% }
%
% @misc{stable_baselines3,
%   author = {A. Raffin and A. Hill and A. Gleave and A. Kanervisto
%             and M. Ernestus and N. Dormann},
%   title  = {Stable-Baselines3: Reliable Reinforcement Learning Implementations},
%   year   = {2021},
%   url    = {https://stable-baselines3.readthedocs.io}
% }
%
% @article{maxpressure2013,
%   author  = {P. Varaiya},
%   title   = {Max pressure control of a network of signalized intersections},
%   journal = {Transportation Research Part C: Emerging Technologies},
%   volume  = {36},
%   pages   = {177--195},
%   year    = {2013},
%   doi     = {10.1016/j.trc.2013.08.014}
% }
%
% @inproceedings{presslight2019,
%   author    = {H. Wei and C. Chen and G. Zheng and K. Wu and
%                V. Gayah and K. Xu and Z. Li},
%   title     = {{PressLight}: Learning Max Pressure Control to Coordinate
%                Traffic Signals in Arterial Network},
%   booktitle = {Proc.\ 25th ACM SIGKDD International Conference on
%                Knowledge Discovery \& Data Mining},
%   pages     = {1290--1298},
%   year      = {2019},
%   doi       = {10.1145/3292500.3330949}
% }
%
% @article{nandakumar2025,
%   author  = {A. Nandakumar and C. Banerjee and L. D. Vanajakshi},
%   title   = {Reinforcement Learning Based Traffic Signal Design to Minimize
%              Queue Lengths},
%   journal = {arXiv preprint arXiv:2509.21745},
%   year    = {2025},
%   url     = {https://arxiv.org/abs/2509.21745}
% }
%
% @article{webster1958,
%   author  = {F. V. Webster},
%   title   = {Traffic Signal Settings},
%   journal = {Road Research Technical Paper},
%   volume  = {39},
%   year    = {1958},
%   note    = {Her Majesty's Stationery Office, London}
% }
%
% @article{mnih2015dqn,
%   author  = {V. Mnih and K. Kavukcuoglu and D. Silver and others},
%   title   = {Human-level control through deep reinforcement learning},
%   journal = {Nature},
%   volume  = {518},
%   pages   = {529--533},
%   year    = {2015},
%   doi     = {10.1038/nature14236}
% }
%
% @article{tan2021hybrid,
%   author  = {T. Tan and W. S. Wong and E. H. Goh and W. H. Chan},
%   title   = {Traffic Signal Control Using Hybrid Action Space Deep
%              Reinforcement Learning},
%   journal = {IEEE Access},
%   volume  = {9},
%   pages   = {52886--52898},
%   year    = {2021},
%   doi     = {10.1109/ACCESS.2021.3070868}
% }